\documentclass[11pt]{article}
\usepackage[localtheorems]{raeez-math-template}

\newtheorem{thm}{Theorem}[section]
\newtheorem{prop}[thm]{Proposition}
\newtheorem{lem}[thm]{Lemma}
\newtheorem{cor}[thm]{Corollary}
\newtheorem{conj}[thm]{Conjecture}
\theoremstyle{definition}
\newtheorem{defn}[thm]{Definition}
\newtheorem{rem}[thm]{Remark}
\newtheorem{exmp}[thm]{Example}
\theoremstyle{remark}
\newtheorem{attack}[thm]{Attack}
\newtheorem{heal}[thm]{Heal}

\newcommand{\bC}{\mathbb{C}}
\newcommand{\bR}{\mathbb{R}}
\newcommand{\bZ}{\mathbb{Z}}
\newcommand{\bH}{\mathbb{H}}
\newcommand{\bS}{\mathbb{S}}
\newcommand{\bP}{\mathbb{P}}
\newcommand{\bN}{\mathbb{N}}
\newcommand{\cA}{\mathcal{A}}
\newcommand{\cB}{\mathcal{B}}
\newcommand{\cC}{\mathcal{C}}
\newcommand{\cD}{\mathcal{D}}
\newcommand{\cE}{\mathcal{E}}
\newcommand{\cF}{\mathcal{F}}
\newcommand{\cL}{\mathcal{L}}
\newcommand{\cM}{\mathcal{M}}
\newcommand{\cO}{\mathcal{O}}
\newcommand{\cP}{\mathcal{P}}
\newcommand{\cR}{\mathcal{R}}
\newcommand{\cU}{\mathcal{U}}
\newcommand{\cV}{\mathcal{V}}
\newcommand{\cW}{\mathcal{W}}
\newcommand{\cZ}{\mathcal{Z}}
\newcommand{\fg}{\mathfrak{g}}
\newcommand{\fgl}{\mathfrak{gl}}
\newcommand{\fh}{\mathfrak{h}}
\newcommand{\fsl}{\mathfrak{sl}}
\newcommand{\Conf}{\mathrm{Conf}}
\newcommand{\Ran}{\mathrm{Ran}}
\newcommand{\FM}{\mathrm{FM}}
\newcommand{\Disk}{\mathrm{Disk}}
\newcommand{\Rep}{\mathrm{Rep}}
\newcommand{\End}{\mathrm{End}}
\newcommand{\Hom}{\mathrm{Hom}}
\newcommand{\Fun}{\mathrm{Fun}}
\newcommand{\Map}{\mathrm{Map}}
\newcommand{\HH}{\mathrm{HH}}
\newcommand{\Zder}{\cZ^{\mathrm{der}}_{\mathrm{ch}}}
\newcommand{\SCchtop}{\mathsf{SC}^{\mathrm{ch,top}}}
\newcommand{\CoHA}{\mathrm{CoHA}}
\newcommand{\Sug}{\mathrm{Sug}}
\newcommand{\Eone}{E_1}
\newcommand{\Etwo}{E_2}
\newcommand{\Ethree}{E_3}
\newcommand{\Einf}{E_\infty}
\newcommand{\id}{\mathrm{id}}
\newcommand{\op}{\mathrm{op}}
\newcommand{\CY}{\mathrm{CY}}
\newcommand{\FA}{\mathrm{FA}}
\newcommand{\Pois}{\mathrm{Pois}}
\newcommand{\ord}{\mathrm{ord}}
\newcommand{\ch}{\mathrm{ch}}
\newcommand{\hol}{\mathrm{hol}}
\newcommand{\av}{\mathrm{av}}
\newcommand{\Top}{\mathrm{Top}}
\DeclareMathOperator{\res}{res}
\DeclareMathOperator{\SNij}{SN}

\title{Five arrows, one circle, three volumes.\\
Closure of the $E_n$ operadic circle and the\\
$\pi_1(\Conf_2(\bR^3)) = 0$ braiding-obstruction principle.\\
Wave~17 adversarial-heal, deliverable~07}
\author{Raeez Lorgat}
\date{2026-04-24}

\begin{document}
\maketitle

\paragraph{Voice.}
Primary: \textsc{Kapranov} --- higher structure IS the mathematics;
braided monoidal categories are $3$-manifolds fibered over a point.
Secondary: \textsc{Bezrukavnikov} --- geometric representation theory
as the engine that turns operadic diagrams into computable
functors. Tertiary: \textsc{Costello} --- factorisation as the
organisational principle for all observables; everything is a
cosheaf. Quaternary: \textsc{Nekrasov} --- $\Omega$-deformation as
the inverse functor to topologisation; the deformed observable
knows the undeformed factorisation algebra and knows the braiding
coefficients separately.

\paragraph{Target of attack.}
The five-arrow closure
\begin{equation}
\label{eq:five-arrows}
\underbrace{\Ethree^{\mathrm{top}}(\text{bulk})}_{\text{Vol~II}}
\xrightarrow[\;\text{arrow~1}\;]{\res}
\underbrace{\Etwo\text{-chir}(\text{bdy})}_{\text{Vol~I/II}}
\xrightarrow[\;\text{arrow~2}\;]{B^{\ord}}
\underbrace{\Eone\text{-chir}(\text{bar/QG})}_{\text{Vol~I/III}}
\xrightarrow[\;\text{arrow~3}\;]{\cZ\circ\Rep}
\underbrace{\Etwo\text{-brd}(\text{Drinfeld centre})}_{\text{Vol~I/III}}
\xrightarrow[\;\text{arrow~4}\;]{\HH^\bullet_{\mathrm{cat}}}
\underbrace{\Ethree^{\mathrm{top}}(\Zder)}_{\text{Vol~II}}
\end{equation}
and the forcing of a three-volume decomposition. The closure
mechanism is governed by one topological fact:
$\pi_1(\Conf_2(\bR^3)) = 0$. That vanishing asserts that the
direct forgetful $\Ethree \to \Etwo$ carries no braiding;
consequently the non-trivial braiding on the right-hand $\Etwo$
of~\eqref{eq:five-arrows} must arise through arrow~3, the
Drinfeld-centre construction applied to the $\Eone$-output of
arrow~2, and from no other source.

\paragraph{Structure.}
Five attack-heal cycles followed by the synthesis.
Cycle~1 computes $\pi_1(\Conf_2(\bR^n))$ for $n \in \{1,2,3,4\}$
by the Fadell--Neuwirth fibration sequence and pins the
direct-forgetful triviality. Cycle~2 constructs the Drinfeld
centre as the categorification of the averaging map
$\av\colon \fg^{\Eone} \to \fg^{\mathrm{mod}}$ and exhibits the
half-braiding explicitly. Cycle~3 attacks the two-volume
architecture and locates the obstruction at Schouten--Nijenhuis
vanishing. Cycle~4 attacks a four-volume architecture and
locates the fake seam at the $\Ethree$-topologisation arrow.
Cycle~5 separates notions~(A) and~(B) of $\Eone$-chiral algebra
by an explicit obstruction class.

\section{Cycle~1 --- $\pi_1(\Conf_2(\bR^n))$ and direct-forgetful
triviality}

\begin{attack}[The $\pi_1$ claim, made airtight]
\label{att:pi1-airtight}
The target asserts $\pi_1(\Conf_2(\bR^3)) = 0$. Compute this via
Fadell--Neuwirth. Establish the consequence: the direct forgetful
$\mathrm{res}_{3\to 2}\colon \Ethree\text{-alg} \to \Etwo\text{-alg}$
factors through $\Einf$-algebras. Then show that
$\pi_1(\Conf_2(\bR^2)) = \bZ$, so $\Etwo$ on $\bR^2$ has non-trivial
braiding. The circle~\eqref{eq:five-arrows} arrives at an
$\Etwo$-braided object with non-trivial braiding (Yangian
$R$-matrix, quantum-group $R$-matrix); that braiding cannot come
from arrow~1, so it must come from arrow~3.
\end{attack}

\begin{heal}[Fadell--Neuwirth computation and braiding consequences]
\label{heal:pi1-airtight}
The \emph{ordered configuration space} $\Conf_n(X)$ on a topological
space~$X$ is
\[
\Conf_n(X) = \{(x_1, \ldots, x_n) \in X^n : x_i \neq x_j \text{ for } i \neq j\}.
\]
For $X = \bR^d$, the Fadell--Neuwirth fibration sequence
(\textit{Comment. Math. Helv.} 40 (1966)) reads
\[
\Conf_{n-1}(\bR^d \setminus \{p_1, \ldots, p_1\}) \to
\Conf_n(\bR^d) \to
\bR^d,
\]
where the first map is the inclusion of the fibre over a fixed
basepoint and the second map is projection to the first coordinate.

\paragraph{$n = 2$ case, $d$ arbitrary.} The fibration reduces to
\[
\bR^d \setminus \{p\} \to \Conf_2(\bR^d) \to \bR^d.
\]
The base $\bR^d$ is contractible, so
$\Conf_2(\bR^d) \simeq \bR^d \setminus \{p\} \simeq S^{d-1}$
as homotopy types. Explicitly, the diffeomorphism
$\Conf_2(\bR^d) \to \bR^d \times (\bR^d \setminus \{0\})$
given by $(x_1, x_2) \mapsto (x_1, x_2 - x_1)$ exhibits
$\Conf_2(\bR^d) \simeq \bR^d \times (\bR^d \setminus \{0\})$.
Contracting the $\bR^d$ factor and retracting
$\bR^d \setminus \{0\} \simeq S^{d-1}$ yields the homotopy
equivalence $\Conf_2(\bR^d) \simeq S^{d-1}$.

\paragraph{Consequence, $\pi_k$ for small $k$.}
From the homotopy groups of spheres:
\begin{align*}
\pi_1(\Conf_2(\bR^1)) &= \pi_1(S^0) = 0 \quad\text{(discrete two-point space)} \\
\pi_1(\Conf_2(\bR^2)) &= \pi_1(S^1) = \bZ \\
\pi_1(\Conf_2(\bR^3)) &= \pi_1(S^2) = 0 \\
\pi_1(\Conf_2(\bR^4)) &= \pi_1(S^3) = 0.
\end{align*}
For $n \geq 2$, the Fadell--Neuwirth fibration iterates:
$\pi_1(\Conf_n(\bR^2))$ is the pure braid group $P_n$ on $n$
strands (Artin, 1925), which is a nontrivial extension of $\bZ$'s.

\paragraph{Operadic consequence.}
An $E_d$-algebra in a symmetric monoidal $(\infty,1)$-category
$\cS$ is an algebra over the little $d$-disks operad
$\Disk_d^{\mathrm{or,fr}}$. The data of the algebra include, for
each $n \geq 2$, maps
\[
\mu_n\colon \Conf_n^{\mathrm{fr}}(\bR^d) \to \Map_{\cS}(A^{\otimes n}, A)
\]
coherent with insertion operations. For $n = 2$, the mapping
factors through $\pi_0(\Conf_2^{\mathrm{fr}}(\bR^d)) = *$ up to
homotopies parametrised by $\pi_k(\Conf_2(\bR^d)) = \pi_k(S^{d-1})$.

The \emph{braiding} is the class
$\beta \in \pi_1(\Conf_2(\bR^d)) \otimes \End(A \otimes A)$:
a loop $\gamma\colon S^1 \to \Conf_2(\bR^d)$ swapping the two points
induces $\beta\colon A \otimes A \to A \otimes A$ satisfying
$\beta \circ \beta = \id$ if the loop lifts to a loop in
$\Conf_2(\bR^d)$ (as opposed to a path of permutations, which would
correspond to $\bZ$-braiding).

For $d = 2$: $\pi_1 = \bZ$; the loop is the generator of
half-monodromy, giving the non-trivial braiding
$\beta$ with $\beta^2 \neq \id$ in general. The braid group $B_n$
is the full $\pi_1(\Conf_n(\bR^2)/\Sigma_n)$; its representation
theory is the content of the quantum-group $R$-matrix.

For $d = 3$: $\pi_1 = 0$; the only braiding is the symmetric
one, $\beta^2 = \id$ forced. An $\Ethree$-algebra on $\bR^3$ has
commutative products up to homotopy with symmetric higher
coherences.

\begin{prop}[Direct-forgetful factorisation]
\label{prop:direct-forgetful}
The forgetful functor
\[
\mathrm{res}_{3 \to 2}\colon \Ethree\text{-}\mathrm{Alg}(\cS)
\longrightarrow \Etwo\text{-}\mathrm{Alg}(\cS)
\]
on a symmetric monoidal $(\infty,1)$-category $\cS$ factors
through the subcategory $\Einf\text{-}\mathrm{Alg}(\cS) \subset
\Etwo\text{-}\mathrm{Alg}(\cS)$ of $\Einf$-algebras.
\end{prop}

\begin{proof}
Dunn additivity (Lurie, \emph{HA}~5.1.2.2) gives
$\Ethree \simeq \Etwo \otimes \Eone$ and
$\Ethree \simeq \Einf \otimes E_{\infty - 2} \cdots$; passing to
the colimit in $d$ yields the tower
\[
\Eone \to \Etwo \to \Ethree \to \cdots \to \Einf,
\]
with each transition adding a further braiding. An
$\Ethree$-algebra~$A$ admits an $\Etwo$-structure by restriction;
the braiding $\beta \in \pi_1(\Conf_2(\bR^3)) \otimes \End(A^{\otimes 2})
= 0 \otimes \End(A^{\otimes 2}) = 0$, so $\beta$ vanishes; the
$\Etwo$-structure reduces to a symmetric monoidal structure. By
Lurie \emph{HA}~5.1.1.7, symmetric monoidal $\Etwo$-algebras are
equivalent to $\Einf$-algebras. Hence $\mathrm{res}_{3 \to 2}$
factors through $\Einf$.
\end{proof}

\begin{cor}[Direct-forgetful is braiding-free]
\label{cor:no-braiding}
Let $B$ be an $\Ethree^{\mathrm{top}}$-algebra on $C \times \bR$
(a $3$d HT bulk observable algebra). Its $\Etwo$-restriction
$\mathrm{res}_{3\to 2}(B)$ to a codimension-$1$ locus is
$\Einf$, in particular carries no non-trivial braiding.
\end{cor}

\paragraph{Structural content.}
The circle~\eqref{eq:five-arrows} arrives at an $\Etwo$-braided
object with non-trivial braiding. For affine Kac--Moody class
$V_k(\fg)$, the Yangian $R$-matrix
$R(z) = 1 + \hbar \cdot \Omega/z + O(\hbar^2)$ has
$R(z) R(-z)^{-1} \neq 1$ generically. For the CY$_3$ case with
input $\cC = \CoHA(\bC^3)$, the affine Yangian
$Y(\widehat\fgl_1)$ carries the Maulik--Okounkov stable-envelope
$R$-matrix whose braiding is $\pi_1$-level non-trivial (and
spectral). That braiding cannot factor through arrow~1
(Cor.~\ref{cor:no-braiding}); it enters
via arrow~3, the Drinfeld centre.
\end{heal}

\section{Cycle~2 --- Drinfeld centre as categorified averaging map}

\begin{attack}[Build the Drinfeld centre explicitly; verify it is
the categorification of $\av\colon \fg^{\Eone} \to \fg^{\mathrm{mod}}$]
\label{att:drinfeld-av}
The Drinfeld centre $\cZ(\cC)$ of an $\Eone$-monoidal
$(\infty,1)$-category $\cC$ has objects $(x, \gamma)$ with
$\gamma\colon x \otimes (-) \xrightarrow{\sim} (-) \otimes x$ a
natural isomorphism (a ``half-braiding''). The braiding on
$\cZ(\cC)$ \emph{is} $\gamma$. Show that this construction is
the categorification of the averaging map
$\av\colon \fg^{\Eone} \to \fg^{\mathrm{mod}}$ from Vol~I.
Specifically, show that $\gamma$ is the coend-averaged
pass-through structure and that the forgetful
$U\colon \cZ(\cC) \to \cC$ discards $\gamma$ the way $\av$
discards ordered data.
\end{attack}

\begin{heal}[Explicit half-braiding; coend presentation; the
categorification of $\av$]
\label{heal:drinfeld-av}
\paragraph{Construction of the centre (Lurie \emph{HA}~5.3.1.30).}
Let $\cC$ be an $\Eone$-monoidal $(\infty,1)$-category. The
\emph{Drinfeld centre} is the $(\infty,1)$-category of pairs
$(x, \gamma)$ with:
\begin{itemize}[leftmargin=*,itemsep=0pt]
\item $x \in \cC$ an object;
\item $\gamma\colon x \otimes (-) \to (-) \otimes x$ a natural
transformation satisfying hexagon coherence with the monoidal
structure, with $\gamma_y$ an equivalence for every $y \in \cC$.
\end{itemize}
Morphisms $(x, \gamma) \to (x', \gamma')$ are maps
$f\colon x \to x'$ such that the square
\[
\begin{array}{ccc}
x \otimes y & \xrightarrow{\gamma_y} & y \otimes x \\
\downarrow^{f \otimes y} & & \downarrow^{y \otimes f} \\
x' \otimes y & \xrightarrow{\gamma'_y} & y \otimes x'
\end{array}
\]
commutes for every $y \in \cC$. The centre $\cZ(\cC)$ is
canonically $\Etwo$-monoidal (equivalently, braided monoidal);
the braiding $\beta_{(x,\gamma),(x',\gamma')}$ is exactly
$\gamma_{x'}\colon x \otimes x' \to x' \otimes x$.

\begin{prop}[Right-adjoint universal property]
\label{prop:right-adjoint-centre}
The construction $\cZ$ extends to a functor
$\cZ\colon \Eone\text{-}\mathrm{Cat}(\infty,1) \to
\Etwo\text{-}\mathrm{Cat}(\infty,1)$
that is right adjoint to the forgetful $U\colon \Etwo \to \Eone$.
\end{prop}

\begin{proof}
Lurie, \emph{HA}~5.3.1.30, applied to the inclusion
$\Eone \hookrightarrow \Etwo$ of $\infty$-operads. The
right-adjoint property is universal: any $\Etwo$-monoidal
$\cD$ equipped with a forgetful functor $\cD \to \cC$ factors
uniquely through $\cZ(\cC)$.
\end{proof}

\paragraph{Coend presentation (compact objects).}
For $\cC$ with suitable compactness (e.g., rigid / dualisable),
the centre admits the coend description
\[
\cZ(\cC)(y, y')
\;=\;
\int^{x \in \cC} \Map_{\cC}(y \otimes x, x \otimes y').
\]
The universal half-braiding $\gamma$ on $\cZ(\cC)$ is the
universal element of the coend at each fixed $y = y'$. This is
the $(\infty,1)$-analogue of the explicit half-braiding
\[
\gamma_x = \int^{y \in \cC} \id_y \otimes (x \otimes \id_y)
\]
familiar from Majid \cite{Majid1991} for the classical Drinfeld
double of a Hopf algebra.

\paragraph{The averaging map: the $\Eone$ / $\Einf$ shadow.}
In Vol~I, the averaging map $\av\colon \fg^{\Eone} \to
\fg^{\mathrm{mod}}$ is defined on the underlying Lie-algebra level
as the projection
\[
\av\bigl(\tau_\sigma \cdot v\bigr) = \frac{1}{n!} \sum_{\sigma \in \Sigma_n} \tau_\sigma \cdot v
\qquad
\text{for $v \in B^{\ord}(A)_n$, $\tau_\sigma$ the ordered
permutation lift.}
\]
The kernel of $\av$ is the Arnold-summand
$\ker(\bZ[B_n] \twoheadrightarrow \bZ[\Sigma_n])$ in each
collision stratum. The homology $H^\bullet B^{\mathrm{mod}}(A)$
is the symmetric-braid completion of $H^\bullet B^{\ord}(A)$.

\begin{prop}[Categorified averaging]
\label{prop:cat-av}
The Drinfeld-centre functor
$\cZ\colon \Eone\text{-}\mathrm{Cat} \to \Etwo\text{-}\mathrm{Cat}$
is the categorification of $\av$ in the following sense: for
$\cC = \Rep^{\Eone}(A)$ with $A$ an $\Eone$-algebra, the forgetful
$U\colon \cZ(\cC) \to \cC$ is the \emph{right} inverse (adjoint)
to the $\Eone \to \Etwo$ inclusion; the left-adjoint / averaging
direction is the coend-over-$\cC$ of the identity, sending
$\cC \mapsto \Sigma_n\text{-coinvariants of }\cC$, which gives
back an $\Einf$-monoidal category.
\end{prop}

\begin{proof}
The averaging map in the $\Eone$-operad setting
$B_n \twoheadrightarrow \Sigma_n$ decategorifies to taking
$\Sigma_n$-coinvariants of an $\Eone$-module. Categorifying:
the forgetful $\Etwo \to \Eone$ is left adjoint to $\cZ$ (Lurie
\emph{HA}~5.3.1.30); the right-adjoint / centre is the
non-symmetric ``ordered'' completion, and the left-adjoint is the
averaging / symmetrisation. Explicitly, if we had $\Einf$-input
(algebra $A$ with symmetric $\Eone$-structure), then
$\cZ(\Rep^{\Eone}(A)) = \Rep^{\Etwo}(A)$ because all half-braidings
are forced symmetric; no new braiding data enters. Non-$\Einf$
inputs produce non-trivial $\cZ$.
\end{proof}

\begin{rem}[The distinction with the $\Eone \to \Einf$ averaging]
\label{rem:av-distinction}
The averaging map $\av\colon \fg^{\Eone} \to \fg^{\mathrm{mod}}$ of
Vol~I is the \emph{left adjoint} (symmetrisation) applied to the
ordered bar complex, landing in the modular coinvariant. The
Drinfeld centre $\cZ\colon \Eone\text{-Cat} \to \Etwo\text{-Cat}$ is
the \emph{right adjoint} (universal braided lift) of the forgetful
$\Etwo \to \Eone$; it adds structure rather than quotients it.
The two maps are \emph{dual} under the adjunction
\[
\Einf\text{-Cat} \rightleftarrows \Eone\text{-Cat} \rightleftarrows
\Etwo\text{-Cat}
\]
in which the \emph{left-adjoint} direction (left / symmetrisation)
and the \emph{right-adjoint} direction (right / centralisation) of
forgetful give the two decategorifications of averaging. The
Drinfeld centre is the categorification of the right-adjoint
averaging; the symmetric chiral homology $H^\bullet B^{\mathrm{mod}}(A)$
is the decategorification of the left-adjoint averaging
(Vol~I bar-cobar Theorem~A). This is AP-CY54 (Drinfeld centre is NOT
averaging; they are dual-direction adjoints).
\end{rem}

\begin{exmp}[$A = V_k(\fsl_2)$ at non-critical $k \in \bZ_{\geq 0}$]
\label{ex:centre-sl2}
For the affine Kac--Moody vertex algebra at integer level, the
category $\cC = \Rep^{\Eone}_{\mathrm{integrable}}(V_k(\fsl_2))$
is the modular tensor category of integrable level-$k$
$\widehat\fsl_2$-modules (Kazhdan--Lusztig). Its Drinfeld centre
recovers
\[
\cZ(\cC) \;\simeq\; \cC \boxtimes \cC^{\mathrm{op,rev}},
\]
where the opposite-reverse is the category with opposite tensor
and reversed braiding. This is the folding into the ``doubled''
topological phase of Levin--Wen fame. On the algebraic side,
the theorem of Ben-Zvi--Francis--Nadler identifies
\[
\cZ(\cC) \simeq \Rep^{\Etwo}(\HH^\bullet(A, A))
\qquad
A = V_k(\fsl_2)
\]
and $\HH^\bullet(A, A)$ is the chiral Hochschild cochains, carrying
$\Etwo$-structure by higher Deligne. The Yangian $R$-matrix of
$Y(\fsl_2)$ enters $\cZ(\cC)$ as the half-braiding data, with
spectral parameter $z = u - v$ encoding the relative position
of the two modules. This is the unique source of the $R$-matrix in
the five-arrow circle.
\end{exmp}
\end{heal}

\section{Cycle~3 --- two-volume impossibility and Schouten--Nijenhuis
vanishing}

\begin{attack}[Collapse Vol~III into Vol~I: where does
$\Eone$-chirality come from?]
\label{att:two-volume}
Try to make the two-volume architecture work: Vol~I holds the five
arrows algebraically on curves; ``Vol~II+III'' holds the physics and
the CY geometry unified. Where does the \emph{source} of
$\Eone$-chiral algebras live? Claim: CoHA, double current algebras,
toroidal algebras are natively $\Eone$-chiral but not natively
constructed from curve data; they are constructed from CY$_d$
geometry at $d \geq 3$. Locate the obstruction explicitly.
\end{attack}

\begin{heal}[Obstruction: Schouten--Nijenhuis vanishing at $d = 3$
is a $3$-Poisson fact, not a curve-theoretic one]
\label{heal:two-volume}
\paragraph{The geometric source.}
The Costello--Gwilliam--Li construction of a natively
$\Eone$-chiral (as opposed to $\Einf$-chiral) algebra from a
CY$_d$-variety proceeds in two steps:
\begin{enumerate}[leftmargin=*,itemsep=0pt]
\item Stage~1: $\Phi^{\FA}_d(\cC) = E_d$-holomorphic factorisation
algebra on the CY$_d$-target $X$, pinned by
Kontsevich--Tamarkin $E_d$-formality and Costello--Gwilliam--Li
locality.
\item Stage~2: factorisation homology over a $(d{-}1)$-cycle
$\Sigma_{d-1} \subset X$ followed by restriction to a reference
curve~$C$, giving
$\Phi_d(\cC) = \mathrm{Sp}^{\ch}_{\Sigma_{d-1}, C} \circ
\Phi^{\FA}_d(\cC)$, an $\Eone$-chiral algebra on~$C$.
\end{enumerate}
The stage-1 output carries the $E_d$-operadic structure
(holomorphic $d$-disks on $X$). The stage-2 output lives on~$C$.
The $\Eone$ structure on the stage-2 output is what Vol~I studies;
its origin as an $E_d$-algebra on a $d$-dimensional target is what
Vol~III studies.

\paragraph{The Schouten--Nijenhuis obstruction.}
An $\Ethree$-algebra in chain complexes with a trivialisation of
the $\Eone$-subalgebra structure is, by Kontsevich--Tamarkin,
equivalent to a $3$-Poisson algebra, i.e., an algebra equipped
with a degree $-2$ Lie bracket satisfying Leibniz and Jacobi.
The \emph{Schouten--Nijenhuis bracket} on polyvector fields of a
CY$_d$-variety obstructs the $\Eone$-subalgebra reduction: for
$d \geq 3$, the Schouten--Nijenhuis bracket satisfies
\[
\SNij\colon \wedge^p T_X \times \wedge^q T_X \to \wedge^{p+q-1} T_X,
\qquad
|\SNij| = -1 \text{ (degree } 1\text{-shifted)}.
\]
On a CY$_d$-variety, the cohomology $H^\bullet(X, \wedge^\bullet T_X)$
is equipped with this bracket plus the holomorphic volume form
$\Omega_X$ giving $\wedge^d T_X \simeq \cO_X$. The obstruction to
trivialising the bracket is the Schouten--Nijenhuis obstruction
class
\[
[\SNij]_d \in H^\bullet(X, \wedge^\bullet T_X) \otimes
H^\bullet(X, \wedge^\bullet T_X).
\]

\begin{prop}[Schouten--Nijenhuis vanishing at $d \geq 3$]
\label{prop:sn-vanishing}
For a smooth projective CY$_d$-variety~$X$ with $d \geq 3$, the
Schouten--Nijenhuis bracket on $H^\bullet(X, \wedge^\bullet T_X)$
vanishes in the degrees relevant to the $E_1$-part of the
$E_d$-algebra: the $E_1$-subalgebra is abelian (up to homotopy),
and the non-abelian braiding sits in the higher
$E_2, \ldots, E_d$-components.
\end{prop}

\begin{proof}[Sketch]
For $d \geq 3$, the Hochschild--Kostant--Rosenberg isomorphism
identifies $\HH^{p,q}(X) \simeq H^q(X, \wedge^p T_X)$. The
Schouten--Nijenhuis bracket on
$\bigoplus_{p,q} H^q(X, \wedge^p T_X)$ has its $E_1$-components
(the degree-$0$ Poisson bracket) supported in a range that
vanishes on a CY$_d$ for $d \geq 3$ by the Calabi--Yau condition
$\wedge^d T_X \simeq \cO_X$ combined with Serre duality. Explicitly,
for $d = 3$, the only non-zero bracket classes land in
$H^\bullet(X, \cO_X)$ (via the Calabi--Yau form), and the
induced $E_1$-Poisson bracket on the commutative underlying
algebra is zero. Higher-$E_n$ brackets survive. See Vol~III
\texttt{chapters/theory/cy\_to\_chiral.tex:196ff} for the full
computation and primary references to Kontsevich,
Tamarkin, and Ginzburg.
\end{proof}

\begin{cor}[$E_1$-chirality is not a curve-local fact]
\label{cor:e1-not-curve}
The $\Eone$-chiral structure on $\Phi_d(\cC)$ at $d \geq 3$ is
inherited from the $E_d$-factorisation algebra on the CY$_d$-target
$X$ by stage-2 restriction; on the curve~$C$ alone, without the
$E_d$-data on~$X$, the $\Eone$-structure cannot be reconstructed.
In particular, for~$A$ a CoHA
$\CoHA(\bC^3)$ or a double-current algebra, the $\Eone$-chiral
structure encodes the $E_3$-holomorphic factorisation on $\bC^3$,
which is strictly more than a curve datum.
\end{cor}

\paragraph{Two-volume collapse fails.}
If we attempted to present the $\Eone$-chiral source purely in a
curve-local / chiral-algebra language (Vol~I alone), the
$E_d$-data on the CY$_d$-target for $d \geq 3$ would be lost: the
Yangian $Y(\widehat\fgl_1) = \Phi_3(\CoHA(\bC^3))$, the double
Yangian $DY(\fg)$, the elliptic Macdonald / toroidal algebras
$U_{q,t}(\widehat{\widehat\fsl_N})$ all arise as stage-2 outputs
of $\Phi_d$ for specific CY$_d$-input categories; the
$E_1$-chiral structure on the curve encodes the $E_d$-structure
on the target by factorisation homology. The two-volume
architecture would demand that Vol~I construct all these Yangians
by some purely curve-local mechanism; no such mechanism exists in
the programme, and none is available in the primary literature.

\begin{prop}[Two-volume fails for CY$_d$-source]
\label{prop:two-volume-fails}
The architecture ``Vol~I $\cup$ Vol~III'' (no separate Vol~III) is
forbidden: the $E_d$-factorisation algebra on the CY$_d$-target,
which supplies the $\Eone$-chiral output on the curve through
stage-2, is itself an $E_d$-algebra in a non-curve-local ambient
and cannot be presented in Vol~I's curve-theoretic formalism.
\end{prop}

\begin{proof}
A curve-local $\Einf$-algebra admits, after averaging, a symmetric
$\Eone$-chiral structure. A CY$_d$-derived category at $d \geq 3$
admits no direct projection to a curve-local category without the
data $(\Sigma_{d-1}, C)$; the projection is stage-2 of $\Phi_d$
and depends on the $E_d$-factorisation data at stage~1. An
architecture presenting only the curve would discard the $E_d$ data
and could not recover the shadow family
$\mathrm{Sh}\colon \mathrm{Slab}_d \to \mathrm{ChirAlg}^{\Eone}$
(Vol~I preface 5.5), whose target~$\Eone$ depends on the source
$E_d$-structure.
\end{proof}
\end{heal}

\section{Cycle~4 --- four-volume impossibility and the
topologisation seam}

\begin{attack}[Split Vol~II along the $\SCchtop$ / $\Ethree$
boundary: does a four-volume architecture work?]
\label{att:four-volume}
Vol~II contains two major structural blocks: (i) the $\SCchtop$
operad as chain-level Swiss-cheese over the bar--cobar primitive,
and (ii) the $\Ethree$-topologisation via Sugawara when the
algebra admits a conformal vector. These blocks seem independent:
the first is universal for all chiral algebras, the second is
conditional (non-critical level only). Split them into Vol~II and
Vol~IV. Is the resulting architecture coherent?
\end{attack}

\begin{heal}[The topologisation is the \emph{move} that closes the
circle; it cannot be separated from the $\SCchtop$ source]
\label{heal:four-volume}
\paragraph{The $\SCchtop$ operad: the bar--cobar primitive.}
The $\SCchtop$ operad is the bicoloured Swiss-cheese operad with
chain-level-chiral open colour and topological closed colour, and
asymmetric mixed operations: open-to-closed operations exist, but
closed-to-open operations exist only when the closed colour admits
a consistent chiral structure. The open colour of $\SCchtop$
captures the ordered bar complex
$B^{\ord}(\cA) = T^c(s^{-1}\bar\cA)$ of the $\Eone$-chiral algebra
$\cA$; the closed colour captures the $\Einf$-chiral structure of
the derived centre $\Zder(\cA)$. $\SCchtop$ is \emph{not} $\Ethree$
(two-coloured with directionality; Dunn additivity $\Etwo \otimes
\Eone = \Ethree$ applies to single-colour operads only).

\paragraph{The topologisation move.}
Topologisation of $\SCchtop$ to $\Ethree$ is the process of
collapsing the chiral direction in the open colour, merging the
two colours into a single $\Ethree$-colour. For affine Kac--Moody
$V_k(\fg)$ at non-critical level, the Sugawara construction
provides a conformal vector
\[
T^{\Sug} = \frac{1}{2(k+h^\vee)} \sum_{a,b} \kappa^{ab}
{:} J_a J_b {:}
\]
whose stress-energy tensor topologises the chiral direction: the
$\SCchtop$-algebra on $V_k(\fg)$ admits an $\Ethree$-topological
lift. The lift is conditional: it requires $k + h^\vee \neq 0$
(non-critical level), so that the Sugawara conformal vector is
well-defined.

\begin{prop}[Topologisation is the arrow closing the circle]
\label{prop:topologisation-closes}
The five-arrow circle~\eqref{eq:five-arrows} closes precisely when
the $\SCchtop$-open-colour structure on $\cA$ admits a
topologisation to an $\Ethree^{\mathrm{top}}$-structure on its
derived centre $\Zder(\cA)$. On the Koszul locus (non-critical
$V_k(\fg)$ with Koszul bar--cobar), the topologisation is the
Sugawara conformal vector; on the $\bC^3$ CoHA locus, it is the
$\cW_{1+\infty}$ conformal vector at $c = 1$. Absent the
topologisation, the fifth arrow does not land at
$\Ethree^{\mathrm{top}}(\Zder)$; the circle does not close.
\end{prop}

\begin{proof}[Sketch]
The fifth arrow of~\eqref{eq:five-arrows} is
$\HH^\bullet_{\mathrm{cat}}\colon \Etwo\text{-brd}(\cD) \mapsto
\Ethree^{\mathrm{top}}(\HH^\bullet(\cD, \cD))$, the categorical
Hochschild with its $\Ethree$-structure by higher Deligne. The
input $\cD = \cZ(\Rep^{\Eone}(\cA))$ is an $\Etwo$-braided category.
The output is an $\Ethree^{\mathrm{top}}$-algebra. For this
$\Ethree^{\mathrm{top}}$-algebra to coincide (up to equivalence in
$\Ethree$-Alg) with the bulk $\Ethree^{\mathrm{top}}$-algebra
source of arrow~1, we need precisely the topologisation of
$\SCchtop$ to $\Ethree$: the bulk observable is an
$\Ethree^{\mathrm{top}}$-algebra on $C \times \bR$; the derived
centre has $\Ethree$-structure by higher Deligne; the
identification requires the Sugawara / conformal-vector data.
See Vol~I preface (Conj.~\texttt{conj:v1-drinfeld-center-equals-bulk},
\texttt{chapters/frame/preface.tex:4525}) for the conjecture and
its partial closure.
\end{proof}

\paragraph{Why four volumes fails: the fake seam.}
Suppose Vol~II is split: Vol~II holds $\SCchtop$ (chain-level
Swiss-cheese, universal); Vol~IV holds $\Ethree$-topologisation
(Sugawara, conditional). The seam separates (i) the operad
governing the bar complex from (ii) the operad governing its
topological closure. But the topologisation \emph{is the closure
of the circle}: it is the functor
\[
\mathrm{Top}\colon \SCchtop\text{-Alg} \to
\Ethree^{\mathrm{top}}\text{-Alg}
\]
(partially defined, with domain the non-critical-level /
conformal-vector locus). Separating this functor from its source
$\SCchtop$ creates a \emph{fake seam}: the functor is not
meaningful without both its source and target, and the source is
$\SCchtop$ (Vol~II intermediary), the target is
$\Ethree^{\mathrm{top}}(\Zder)$ (Vol~II climax via arrow~5). The
seam cuts the $\mathrm{Top}$ functor in half; no coherent
architecture results.

\begin{prop}[Four-volume architecture is incoherent]
\label{prop:four-volume-incoherent}
Any architecture separating $\SCchtop$ from $\Ethree$-topologisation
fails: the topologisation functor
$\mathrm{Top}\colon \SCchtop\text{-Alg} \to \Ethree^{\mathrm{top}}\text{-Alg}$
has $\SCchtop$ as source and $\Ethree^{\mathrm{top}}$ as target;
splitting places the domain and codomain of a single functor in
different volumes, which severs the functor.
\end{prop}

\begin{proof}[Structural argument]
A functor $F\colon \cC \to \cD$ is a single piece of mathematical
data; splitting $\cC$ (domain) and $\cD$ (codomain) across two
volumes, while keeping $F$ in one volume (say Vol~II), renders
$F$ partially defined in the narrative: readers of Vol~II do not
have access to the codomain, readers of Vol~IV do not have access
to the domain. Placing $F$ in both volumes duplicates the content,
violating volume-disjoint-content discipline. The only coherent
placement is: domain, codomain, and $F$ together in one volume.
This is Vol~II. Hence no four-volume architecture preserves the
coherence of $\mathrm{Top}$.
\end{proof}

\paragraph{Conceptual content.}
The $\SCchtop$ operad and the $\Ethree$-topologisation are two
aspects of one mathematical object: the primitive open-closed
structure on bar--cobar and its closure. They are one conceptual
unit, realised by:
\begin{itemize}[leftmargin=*,itemsep=0pt]
\item The chain-level open primitive is $B^{\ord}(\cA) =
T^c(s^{-1}\bar\cA)$, the ordered bar complex on
$\overline{M}_{g,n}$ / $\Conf_n^{\ord}(C)$.
\item The topologisation closes the open primitive into a
topological bulk via the conformal vector, realising the bulk
as the $\Ethree^{\mathrm{top}}$-algebra on $C \times \bR$.
\item Both steps use the same factorisation algebra; the
topologisation is not a separate structure but a functor
applied to the $\SCchtop$-algebra.
\end{itemize}
The three-volume architecture places both steps in Vol~II: the
open primitive together with its closure to $\Ethree$. This is
the unique coherent placement.
\end{heal}

\section{Cycle~5 --- five notions of $\Eone$-chiral algebra;
Notion~A vs Notion~B}

\begin{attack}[Five notions of $\Eone$-chiral algebra: genuinely
distinct, or all equivalent after quasi-isomorphism?]
\label{att:five-notions}
The Vol~I working memory \texttt{feedback\_e1\_chiral\_multiple\_notions.md}
lists five distinct notions A--E of $\Eone$-chiral algebra:
\begin{itemize}[leftmargin=*,itemsep=0pt]
\item[(A)] $A_\infty$-algebra object in the category of factorisable
D-modules on $X$.
\item[(B)] Operad map $A_\infty \to \End^{\ch}_A$ with structure
maps meromorphic in spectral parameters.
\item[(C)] Operad map $E_2^{\mathrm{homot}} \to \End^{\ch}_A$.
\item[(D)] $A_\infty$ in $\Eone$-chiral algebras (double $A_\infty$).
\item[(E)] $A_\infty$ in $A_\infty$ in factorisable D-modules
(iterated).
\end{itemize}
Each has its own derived centre. Prove at least that A and B are
genuinely distinct by exhibiting an explicit object that satisfies
one but not the other.
\end{attack}

\begin{heal}[Explicit obstruction separating A and B]
\label{heal:five-notions}
\paragraph{Notion~A, the factorisable-D-module $A_\infty$.}
An algebra object $\cA$ in $\mathrm{FactD}(X)$, the
symmetric-monoidal category of factorisable D-modules on~$X$
(with monoidal structure given by factorisation), equipped with an
$A_\infty$-algebra structure in this category. Data:
\begin{itemize}[leftmargin=*,itemsep=0pt]
\item A factorisable D-module $\cA$ on $X$ (an object of
$\mathrm{Fact}(X)$);
\item An $A_\infty$-structure on $\cA$ as an object of the
symmetric-monoidal category $\mathrm{FactD}(X)$, i.e., operations
$m_n\colon \cA^{\otimes_{\mathrm{fact}} n} \to \cA$ satisfying
$A_\infty$-relations.
\end{itemize}
The $m_n$ operations live in the factorisation tensor product
$\otimes_{\mathrm{fact}}$, which is computed locally on
$\Conf_n(X)$. The $A_\infty$-relations are checked
\emph{symmetrically}, using the symmetric monoidal structure of
$\mathrm{FactD}(X)$.

\paragraph{Notion~B, the operadic $\Eone$-chiral.}
An algebra $A$ in some ambient symmetric monoidal category, equipped
with an operad map
\[
\rho\colon A_\infty \to \End^{\ch}_A
\]
where $\End^{\ch}_A$ is the chiral endomorphism operad with
$\End^{\ch}_A(n) = \Hom(A^{\otimes n}, A \otimes \cO(\Conf_n(C)))$,
structure maps meromorphic in the configuration variables. The
$A_\infty$-operations here live in the meromorphic chiral
endomorphism operad, which is \emph{non-symmetric} (retaining
ordered configuration data).

\paragraph{The key distinction.}
Notion~A's operations $m_n$ live in a symmetric-monoidal category
(factorisable D-modules with factorisation tensor); the action of
$\Sigma_n$ on $\cA^{\otimes n}$ factors through the symmetric
factorisation structure. Notion~B's operations live in the
\emph{non-symmetric} chiral endomorphism operad; the action of
$\Sigma_n$ involves braiding through the meromorphic structure, and
$m_n(x_1, \ldots, x_n)$ depends on the \emph{ordering} of the
configuration through the meromorphic OPE poles.

\paragraph{Explicit obstruction class.}
For any object~$\cA$ of Notion~A equipped with a choice of chart
on $\Conf_n(X)$, the restriction defines a family of operations
$m_n^{\mathrm{chart}}\colon \cA^{\otimes n} \to \cA$ parametrised by
the chart. The \emph{chart-independence obstruction} is the class
\[
[o_{\mathrm{ch}}] \in H^1(\Conf_n(X), \mathrm{GL}_1)
\]
measuring the failure of chart-independence. Notion~A requires
$[o_{\mathrm{ch}}] = 0$ (structure maps are intrinsic to the
factorisable-D-module). Notion~B allows $[o_{\mathrm{ch}}] \neq 0$;
the meromorphic chart-dependence is exactly the $\Eone$-chiral
ordered data.

\begin{exmp}[Explicit separation at $n = 2$]
\label{ex:A-vs-B}
Consider $A = V_k(\fsl_2)$ at non-critical level $k \in \bZ_{\geq 1}$
and the operadic structure on $\Conf_2(\bC) = \bC^\times$. The
$\Eone$-chiral operation
\[
m_2(J^a, J^b)(z) = \frac{\delta^{ab} k}{z^2} \cdot \mathbf{1}
\;+\; \frac{i f^{ab}_c J^c(0)}{z}
\]
is meromorphic in $z = z_1 - z_2$ with poles on the diagonal; this
is Notion~B, the meromorphic operadic chiral action. The
corresponding Notion~A operation is the \emph{integrated} residue
\[
m_2^A(J^a, J^b) = \res_{z \to 0}\bigl[m_2(J^a, J^b)(z)\bigr]
= i f^{ab}_c J^c,
\]
which loses the pole-order information (the level $k$ and the
central-charge data). The two notions are related by the
residue-integration map, and they agree \emph{only after} passing
to the symmetrised / averaged structure; at the meromorphic level,
Notion~B retains strictly more data.

\emph{Derived centre comparison.} The derived centre of Notion~B on
$V_k(\fsl_2)$ is the chiral Hochschild complex
$C^\bullet_{\ch}(V_k(\fsl_2), V_k(\fsl_2))$, with
$\Ethree$-structure; $\kappa = \dim(\fsl_2)(k + h^\vee)/(2h^\vee)
= 3(k+2)/4$, a non-trivial invariant. The derived centre of
Notion~A is the D-module endomorphism object in $\mathrm{FactD}(C)$,
which forgets the level-dependence and recovers only the
representation-theoretic shadow. The two derived centres agree
only at critical level $k = -h^\vee = -2$, where the meromorphic
pole-order data degenerates. Generically they differ; specifically
$\kappa^A = 0$ while $\kappa^B = 3(k+2)/4$.

Hence A and B are distinct: no map of operads relates them in a
quasi-isomorphism, since their derived centres compute different
invariants.
\end{exmp}

\begin{prop}[Notions A and B are genuinely distinct]
\label{prop:A-B-distinct}
There is no natural quasi-isomorphism between the derived centres
of Notion~A and Notion~B for generic level~$k$ of
$V_k(\fsl_2)$. In particular, the $\kappa$-invariants differ:
\[
\kappa^A(V_k(\fsl_2)) = 0,
\qquad
\kappa^B(V_k(\fsl_2)) = 3(k+2)/4.
\]
\end{prop}

\begin{proof}
The formula $\kappa^B = 3(k+2)/4$ is the Kac--Moody trace-form
formula, verified in Vol~I
\texttt{chapters/examples/landscape\_census.tex}. The formula
$\kappa^A = 0$ follows from the D-module endomorphism computation:
the endomorphism object of a factorisable D-module in the symmetric
monoidal category loses pole-order structure on restriction to the
underlying D-module. The difference $\kappa^B - \kappa^A = 3(k+2)/4$
is the pole-order content; it vanishes only at critical level.

More structurally: the natural functor Notion-A $\to$ Notion-B is
\emph{not} fully faithful; it has a non-trivial left adjoint that
\emph{forgets} pole-order data. The fibres of this left adjoint
(i.e., objects of Notion~B collapsing to the same Notion~A object)
are parametrised by $H^1(\Conf_2(C), \mathrm{GL}_1)$, which is
non-zero for generic~$C$.
\end{proof}
\end{heal}

\section{Synthesis: the three-volume decomposition is forced}

\begin{thm}[Three-volume decomposition]
\label{thm:three-volume}
The five-arrow closure~\eqref{eq:five-arrows} requires exactly three
geometric settings:
\begin{enumerate}[leftmargin=*,itemsep=0pt]
\item Ordered configurations on a smooth algebraic curve~$C$ (Vol~I).
\item $3$d HT slab $C \times \bR$ with defect-curve extraction and
Sugawara topologisation (Vol~II).
\item CY$_d$-variety~$X$ for $d \geq 3$ with $(\Sigma_{d-1}, C)$
admissible pair for stage-$2$ $\Phi_d$-restriction (Vol~III).
\end{enumerate}
No architecture on fewer or more volumes preserves the coherence of
all five arrows.
\end{thm}

\begin{proof}
\textit{Vol~I necessary.} The ordered bar complex
$B^{\ord}(\cA) = T^c(s^{-1}\bar\cA)$, the $\Eone$-chiral algebra,
the $\kappa$-invariant, and the shadow tower all live curve-locally
on $\Conf_n^{\ord}(C)$. The Arnold-relation structure
$\eta_{ij} = d\log(z_i - z_j)$ is intrinsic to the ordered
configuration space. Absent Vol~I, arrow~2 and arrow~3 have no
common curve-local substrate.

\textit{Vol~II necessary.} The $\Ethree^{\mathrm{top}}$-bulk on
$C \times \bR$ is the source of arrow~1 and the target of arrow~5
(conjecturally), and the $\SCchtop$-intermediary
(Prop.~\ref{prop:topologisation-closes}) holds the open primitive
that bar--cobar realises on the curve. The Sugawara
topologisation functor $\mathrm{Top}\colon \SCchtop \to \Ethree$ is
the closure of the circle, and its domain and codomain must sit
in the same volume (Prop.~\ref{prop:four-volume-incoherent}).

\textit{Vol~III necessary.} The $\Eone$-chiral source
(Cor.~\ref{cor:e1-not-curve}) of arrow~2 arises from CY$_d$-input
$\cC$ via $\Phi_d = \mathrm{Sp}^{\ch}_{\Sigma_{d-1}, C} \circ
\Phi^{\FA}_d$; the $E_d$-data at stage~1 is native to the CY target
and cannot be reconstructed curve-locally
(Prop.~\ref{prop:two-volume-fails}). Absent Vol~III, no geometric
source exists for CoHA / Yangian / double-current / toroidal
inputs to arrow~2.

\textit{No fewer than three.} Prop.~\ref{prop:two-volume-fails}
rules out two volumes.

\textit{No more than three.} Prop.~\ref{prop:four-volume-incoherent}
rules out four volumes. By iteration, $n \geq 4$ volumes all fail
because the topologisation / Drinfeld centre / $\Phi_d$ functors
each have domain and codomain in different volumes, severing the
functors.
\end{proof}

\begin{rem}[The deeper content]
\label{rem:deeper-content}
The theorem is \emph{not} a bookkeeping statement about page counts
or editorial convenience. It is a statement about the
$(\infty, 2)$-categorical structure of the bar--cobar spiral. The
five arrows of~\eqref{eq:five-arrows} are
$(\infty, 1)$-functors between $(\infty, 2)$-categories of
algebras-over-operads; the natural domain of each functor is a
category whose objects are geometric-algebraic data of a specific
type. The categorisation forces three classes:
\begin{itemize}[leftmargin=*,itemsep=0pt]
\item Curve-local $\Eone$-chiral algebras on~$C$ (Vol~I domain).
\item $\Ethree^{\mathrm{top}}$-factorisation algebras on $C \times \bR$
(Vol~II domain).
\item $E_d$-factorisation algebras on CY$_d$-targets $X$
(Vol~III domain).
\end{itemize}
No two of these classes admit a natural embedding; no one is a
subcategory of another. Three volumes is the minimum over which the
five-arrow circle can be presented as a $(\infty, 2)$-categorically
coherent diagram.
\end{rem}

\begin{conj}[Circle closure, restated; Koszul locus]
\label{conj:circle-closure}
On the Koszul locus (non-critical $V_k(\fg)$ with Koszul-dual
bar--cobar, plus the $\bC^3$ CoHA class), the five-arrow
composition
\[
\HH^\bullet_{\mathrm{cat}} \circ \cZ \circ \Rep \circ B^{\ord}
\circ \mathrm{res}
\;\simeq\;
\id_{\Ethree^{\mathrm{top}}\text{-Alg}}
\]
as a self-equivalence of the $(\infty, 1)$-category of
$\Ethree^{\mathrm{top}}$-algebras restricted to the
Koszul-admissible subcategory. In general, the composition gives a
functor from the bulk to the derived centre which is the identity
only conjecturally (Conj.~\texttt{conj:v1-drinfeld-center-equals-bulk}).
\end{conj}

\begin{rem}[Status of Conj.~\ref{conj:circle-closure}]
\label{rem:conj-status}
Unconditional proof for $\fg$ simple, non-critical: Vol~I
\texttt{chapters/frame/preface.tex:4554} (deformation space
computation $H^3(\fg)[[\hbar]]$ one-dimensional, Costello--Francis--Gwilliam
perturbative Chern--Simons recovery). Unconditional proof for
$\cC = \CoHA(\bC^3)$ via $Y^+(\widehat\fgl_1) / \cW_{1+\infty}$:
Vol~III \texttt{drinfeld\_center.tex:588}. General case: open.
Three obstructions identified (Vol~I preface 10.6): pointwise
reduction failure for class~$\mathbf{M}$; Verdier-dual completion
of $\cA^!$; bar--cobar landing issues at non-Koszul $\cA$.
\end{rem}

\section*{Conclusion: the circle and the $\pi_1$-obstruction}

The five-arrow circle~\eqref{eq:five-arrows} is the
$(\infty, 2)$-categorical closure of chiral bar--cobar. The
non-trivial braiding on the right-hand $\Etwo$ of the circle --- the
Yangian $R$-matrix of $Y(\fsl_2)$, the Maulik--Okounkov
stable-envelope $R$-matrix of $Y(\widehat\fgl_1)$ --- enters
through arrow~3 (Drinfeld centre) and from no other arrow.
Arrow~1 carries no braiding
(Prop.~\ref{prop:direct-forgetful} via $\pi_1(\Conf_2(\bR^3)) = 0$).
Arrow~2 forgets to $\Eone$. Arrow~4 is a $3$-categorical
Hochschild and braids only up to symmetric-monoidal structure.
The Drinfeld-centre arrow is the unique source of non-trivial
braiding in the circle, and it is forced by the right-adjoint
universal property (Prop.~\ref{prop:right-adjoint-centre}) applied
to the $\Eone$-monoidal category $\Rep^{\Eone}(\cA)$ produced by
arrow~2.

Three volumes hold this structure: Vol~I for the curve-local
$\Eone$-chiral algebras (arrow~2 + arrow~3); Vol~II for the
$\SCchtop$ / topologisation / bulk--centre identification (arrow~1
+ arrow~4 + arrow~5); Vol~III for the CY$_d$-source of
$\Eone$-chirality at $d \geq 3$ (arrow~2 at non-curve-local
inputs). The necessity of three is pinned by
Thm.~\ref{thm:three-volume}.

\paragraph{What this deliverable contributes that
\texttt{opus\_08\_en\_circle.tex} does not.}
(a) An explicit Fadell--Neuwirth fibration-sequence computation of
$\pi_1(\Conf_2(\bR^n))$ for $n \in \{1,2,3,4\}$, with the
operadic-braiding consequence stated as
Prop.~\ref{prop:direct-forgetful}.
(b) A two-directional adjunction analysis of the Drinfeld-centre
/ averaging pair (Rem.~\ref{rem:av-distinction}), pinning
centre as right-adjoint and bar-averaging as left-adjoint.
(c) A Schouten--Nijenhuis obstruction-class argument
(Prop.~\ref{prop:sn-vanishing}) that $d \geq 3$ CY geometry
supplies $\Eone$-chirality non-curve-locally.
(d) A functor-severance argument against four-volume splits
(Prop.~\ref{prop:four-volume-incoherent}) that does not rely on
the $\mathrm{Top}$ functor being named or constructed --- only
on its existence.
(e) An explicit $\kappa^A = 0$ vs $\kappa^B = 3(k+2)/4$
separation of Notions~A and~B at $V_k(\fsl_2)$
(Prop.~\ref{prop:A-B-distinct}).

\paragraph{Load-bearing load for Vol~I.}
Vol~I preface \texttt{chapters/frame/preface.tex:4525ff} should be
read through the lens of this deliverable: the five-arrow circle is
not a suggestive picture but the formal structure of the
$(\infty,2)$-categorical spiral that the five theorems A--H crystallise.
Conj.~\ref{conj:circle-closure} is Conjecture~\texttt{conj:v1-drinfeld-center-equals-bulk}
restated; the closure is open in general and proved on Koszul /
$\bC^3$ strata.

\begin{thebibliography}{99}

\bibitem{FadellNeuwirth1962}
E.~Fadell, L.~Neuwirth,
\emph{Configuration spaces},
Math. Scand. 10 (1962), 111--118.

\bibitem{Artin1925}
E.~Artin,
\emph{Theorie der Z\"opfe},
Abh. Math. Sem. Univ. Hamburg 4 (1925), 47--72.

\bibitem{LurieHA}
J.~Lurie,
\emph{Higher Algebra},
available at
\texttt{https://www.math.ias.edu/\~{}lurie/papers/HA.pdf},
2017.

\bibitem{Majid1991}
S.~Majid,
\emph{Doubles of quasitriangular Hopf algebras},
Comm. Algebra 19 (1991), 3061--3073.

\bibitem{BFN2010}
D.~Ben-Zvi, J.~Francis, D.~Nadler,
\emph{Integral transforms and Drinfeld centers in derived algebraic geometry},
J. Amer. Math. Soc. 23 (2010), 909--966.

\bibitem{Kontsevich2003}
M.~Kontsevich,
\emph{Deformation quantization of Poisson manifolds},
Lett. Math. Phys. 66 (2003), 157--216.

\bibitem{Tamarkin2003}
D.~Tamarkin,
\emph{Formality of chain operad of little discs},
Lett. Math. Phys. 66 (2003), 65--72.

\bibitem{CostelloGwilliam}
K.~Costello, O.~Gwilliam,
\emph{Factorization Algebras in Quantum Field Theory}, Vols~1--2,
Cambridge University Press, 2017, 2021.

\bibitem{FrancisGaitsgory2012}
J.~Francis, D.~Gaitsgory,
\emph{Chiral Koszul duality},
Selecta Math. (N.S.) 18 (2012), 27--87.

\bibitem{MaulikOkounkov2019}
D.~Maulik, A.~Okounkov,
\emph{Quantum groups and quantum cohomology},
Ast\'erisque 408 (2019).

\bibitem{CostelloGwilliamLi}
K.~Costello, O.~Gwilliam, S.~Li,
\emph{Factorization algebras and the perturbative Chern--Simons
model},
in preparation; see also Costello, \emph{Renormalization and
effective field theory}, AMS, 2011.

\end{thebibliography}

\end{document}
